\begin{frame}
\frametitle{Paměťové prostory procesů a odkládání na disk}

Více procesů, každý vlastní paměťový prostor. Vzájemná ochrana
a možnost rozšíření kapacity hlavní paměti o prostor na sekundární -- swap.

{
\centering

\includegraphics[width=0.80\linewidth]{generated/memory/paging/paging-to-disc.pdf}

}

Překlad po jednotlivých bytech by byl drahý, prostor je rozdělěný na~stránky
(zarovnané), typicky 4\,kB, někdy i větší, např. 64\,kB.

\end{frame}

\begin{frame}
\frametitle{Překlad virtuální a fyzické adresy}

Překlad zajišťuje \textbf{Memory Management Unit} (MMU).

Po vyplnění stránkovacích tabulek a nastavení \textbf{Page Directory Base Register}
(PDBR) operačním systémem probíhá překlad na stránky přítomné v paměti většinou automaticky.

Naopak výpadky stránek ošetřují rutiny operačního systému. Pokud je přístup do namapované
oblasti, načtou data z disku, sítě, odkládacího oddílu.

Stějně jako u cache je zde nutné a náročné hledat prostor pro nově potřebné stránky (princip podobný LRU).

\vskip 2mm

{
\centering

\includegraphics[width=0.80\linewidth]{generated/memory/paging/paging-cpu-mm-mem.pdf}

}
\end{frame}

\begin{frame}
\frametitle{Jednoúrovňová stránkovací tabulka}

{
\centering

\includegraphics[width=0.80\linewidth]{generated/memory/paging/paging-1-level-cz.pdf}

}

Řešení je nevýhodné, pro každý, i malý, spuštěný proces je na 32-bitovém
systému pro 4\,kB stránky potřeba alokovat 4\,MB (překlad 20-bitů adresy,
4-byte na 32-bit položku) 

\end{frame}


\begin{frame}
\frametitle{Dvouúrovňové stránkování}

{
\centering

\includegraphics[width=0.7\linewidth]{generated/memory/paging/paging-2-levels-cz.pdf}

}

Adresář stránek pro každý proces, úroveň tabulek s vlastními položkami
se alokuje jen, když je daná oblast virtuálního prostoru použitá.

Za pozornost stojí, že při překladu po 10-ti bitech vychází pole položek
adresáře i tabulky na velikost 4\,kB. Nepotřebuje alokace vyššího řádu
ohrožené fragmentací.

\end{frame}

\begin{frame}
\frametitle{Položky adresářů stránek -- \textbf{Page Directory Entry} (PDE) }

{
\centering

\includegraphics[width=0.90\linewidth]{generated/memory/paging/paging-pde-1.pdf}

}

\vskip 2mm

\begin{itemize}
\item bit 0: Present bit -- další úroveň nebo stránka je přítomná v paměti (1). Pokud není (0) tak jsou data uložena na disku nebo oblast není mapovaná. Někdy také bit označovaný V -- valid bit.
\item bit 1: Read/Write: 1 -- R/W -- zápis povolený; 0 -- read only (RO)
\item bit 2: User/Supervisor: 1 -- přístupné uživateli; 0 -- přístup jen operační systém
\item bit 3: Write-through/Write-back -- nastavení typu přístupu
\item bit 4: Cache disabled/enabled -- důležité pro paměťově mapované periferie, přímý přístup do registrů bez cache
\item bit 5: Accessed -- nastaven MMU při prvním přístupu z procesoru
\end{itemize}

\end{frame}

\begin{frame}
\frametitle{Položky koncových tabulek -- \textbf{Page Table Entry} (PTE)}

{
\centering

\includegraphics[width=0.90\linewidth]{generated/memory/paging/paging-pte.pdf}

}

\vskip 2mm

\begin{itemize}
\item bit 6: Dirty bit, případně někdy Modified – jen nastavený MMU, pokud došlo k zápisu do dané stránky od poslední kontroly a případného nulování příznaku operačním systémem.
\item Ostatní bity mají stejný význam jako v položce adresáře
\end{itemize}

Bit \textbf{Accessed} je klíčový, když operační systém hledá stránky fyzické paměti k uvolnění.
Obvykle počítá podle bitů další, obsáhlejší, statistiku a při průchodu seznamem bity nuluje.
Pokud je nastavený \textbf{Dirty}, tak je při uvolnění fyzciké stránky (PFN) nutné data
sychronizovat zpět do~souboru, odkládacího oddílu atd.

\end{frame}


\begin{frame}
\frametitle{Překlad na 64-bitových architekturách}

{
\centering

\includegraphics[width=0.90\linewidth]{generated/memory/paging/address-trans-4-levels.pdf}

}

Vyžaduje více úrovní. Pro 64-bitů a 4\,kB stránky je teoreticky
potřeba přeložit 52-bitů. Pokud jednotlivé části tabulek mají
zaplnit jednu stránku, tak vychází překlad na 9\,bitů ($512 \times 8 = 4096$)
na úroveň a tedy $ceil(52 / 9) = 6$. Často se ale jedna až dvě úrovně vynechávají, více později.


\end{frame}

\begin{frame}
\frametitle{Zrychlení opakovaného překladu}

Každý překlad představuje několik přístupů do~paměti a i když je urychlený
přístupy přes cache, tak zpomaluje. Zavádí se \textbf{Translation Look-Aside Buffer} (TLB).
Pracuje na shodném principu jako paměťová cache, ale ukládá k virtuální adrese
její překlad na~fyzickou. Opět omezená kapacita, LRU a~nutnost brát v~potaz
při~programování.

\vskip 2mm

{
\centering

\includegraphics[width=0.80\linewidth]{generated/memory/paging/paging-tlb.pdf}

}
\end{frame}
